%%%%%%%%%%%%%%%%%%%%%%%%%%%%%%%%%%%%%%%%%%%%%%%%%%%%%%%%%%%%%%%%%%%%%%%%%%%%%%%%

\pagestyle{empty}
\begin{abstract}
Este traballo estuda o desenvolvemento dun emulador da consola Game Boy orixinal.
Isto supón o deseño e a implementación dun intérprete do conxunto de instrucións do Sharp LR35902 e do resto do hardware que permite o correcto funcionamento da consola,
como a unidade de procesamento de píxeles (PPU) e os controladores de bancos de memoria (MBC).
O obxectivo é crear un emulador multiplataforma e facilmente reusable, que permita ser integrado en sistemas embebidos.
Para isto creouse un núcleo de emulación independente e dúas implementacións de frontend,
unha adicada a usarse nun dispositivo da familia Raspberry Pi Zero e outra adicada para depuración en sesións gŕaficas completas e no navegador.

  \vspace*{25pt}
  \begin{segundoresumo}
    \blindtext % substitúe este comando polo resumo do teu TFG
               % na lingua secundaria do documento (tipicamente: inglés)
  \end{segundoresumo}
\vspace*{25pt}
\begin{multicols}{2}
\begin{description}
\item [\palabraschaveprincipal:] \mbox{} \\[-20pt]
\begin{itemize}
  \item Emulación
  \item Game Boy
  \item Sistemas embebidos
  \item Multiplataforma
  \item Rust
  \item Raylib
  \item Nix
\end{itemize}
\end{description}

\begin{description}
\item [\palabraschavesecundaria:] \mbox{} \\[-20pt]
\begin{itemize}
  \item Emulation
  \item Game Boy
  \item Embedded systems
  \item Cross-platform
  \item Rust
  \item Raylib
  \item Nix
\end{itemize}
\end{description}
\end{multicols}
\end{abstract}
\pagestyle{fancy}

%%%%%%%%%%%%%%%%%%%%%%%%%%%%%%%%%%%%%%%%%%%%%%%%%%%%%%%%%%%%%%%%%%%%%%%%%%%%%%%%
